% GENERATED FILE. Edit canonical lessons, then run npm run generate:books.
% canonical-source-hash: fnv1a64:f8c8b9215c0b8fba
% canonical-lessons: BN-C314-shotru, BN-C314-soinyo, BN-C314-akromon-kora, BN-C314-hamla, BN-C314-dokhol-kora

\chapter{War Comes}
\label{ch:bn-war-comes}
\hlchaptermodality{314}

\begin{chapteropening}
\textbf{By the end of this chapter:} \emph{I can tell of enemies, soldiers, raids and attacks.}

Complete the last lesson of chapter 314: I can tell of enemies, soldiers, raids and attacks.
\end{chapteropening}

\section[\texorpdfstring{\bn{শত্রু} (shôtru)}{শত্রু (shôtru)}]{\bn{শত্রু} (shôtru) — an enemy}

\label{lesson:BN-C314-shotru}



\begin{quote}
Before the new one: say the Bengali for a crown, then the Bengali for to rule.
\end{quote}

\subsection*{You'll want to know: \bn{শত্রু}}
\textbf{\bn{শত্রু}} — \emph{shôtru} — \textquotedblleft{}an enemy\textquotedblright{}.

\begin{quote}
\bn{শত্রুরা} \bn{রাতে} \bn{নদী} \bn{পার} \bn{হল।}
\end{quote}

\begin{quote}
\emph{shôtrurā rāte nôdī pār hôlô.}
\end{quote}

\begin{quote}
The enemy crossed the river by night.
\end{quote}

\begin{cousinweb}
\textbf{\bn{শত্রু}} (\emph{shôtru}) is Sanskrit \emph{śatru}, 'an enemy, a foe, a rival' (Monier-Williams), unchanged.
\end{cousinweb}

\begin{grammarlens}[title={What you've built}]
One more word for telling what happened.
\end{grammarlens}

\subsection*{Your turn}
\begin{itemize}
\raggedright
  \item \emph{Say it:} \emph{shôtru}
  \item \emph{Say it:} \emph{shôtru}, once more
  \item \emph{Say it:} \emph{shôtrurā rāte nôdī pār hôlô.} — The enemy crossed the river by night.
  \item \emph{Recall:} say the Bengali for a crown, then the Bengali for to rule, then say \emph{shôtru} again
\end{itemize}

\begin{tcolorbox}[breakable,colback=teal!4,colframe=teal!35!black,title={Before you move on}]
What does \bn{শত্রু} mean? (\textquotedblleft{}An enemy\textquotedblright{}.) Say it once more.
\end{tcolorbox}
\section[\texorpdfstring{\bn{সৈন্য} (soinyô)}{সৈন্য (soinyô)}]{\bn{সৈন্য} (soinyô) — a soldier}

\label{lesson:BN-C314-soinyo}



\begin{quote}
Before the new one: say the Bengali for to rule, then the Bengali for an enemy.
\end{quote}

\subsection*{You'll want to know: \bn{সৈন্য}}
\textbf{\bn{সৈন্য}} — \emph{soinyô} — \textquotedblleft{}a soldier\textquotedblright{}.

\begin{quote}
\bn{রাজার} \bn{সৈন্যরা} \bn{দুর্গের} \bn{দরজায়} \bn{দাঁড়াল।}
\end{quote}

\begin{quote}
\emph{rājār soinyôrā durger dôrjāy dā̃ṛālô.}
\end{quote}

\begin{quote}
The king's soldiers took up their places at the gate of the fort.
\end{quote}

\begin{cousinweb}
\textbf{\bn{সৈন্য}} (\emph{soinyô}) is Sanskrit \emph{sainya}, 'belonging to an army', and as a noun 'a soldier' (Monier-Williams); its \bn{ৈ} is read \emph{oi}. An army is a \textbf{\bn{সেনাবাহিনী}}.
\end{cousinweb}

\begin{grammarlens}[title={What you've built}]
One more word for telling what happened.
\end{grammarlens}

\subsection*{Your turn}
\begin{itemize}
\raggedright
  \item \emph{Say it:} \emph{soinyô}
  \item \emph{Say it:} \emph{soinyô}, once more
  \item \emph{Say it:} \emph{rājār soinyôrā durger dôrjāy dā̃ṛālô.} — The king's soldiers took up their places at the gate of the fort.
  \item \emph{Recall:} say the Bengali for to rule, then the Bengali for an enemy, then say \emph{soinyô} again
\end{itemize}

\begin{tcolorbox}[breakable,colback=teal!4,colframe=teal!35!black,title={Before you move on}]
What does \bn{সৈন্য} mean? (\textquotedblleft{}A soldier\textquotedblright{}.) Say it once more.
\end{tcolorbox}
\section[\texorpdfstring{\bn{আক্রমণ} \bn{করা} (ākrômôn kôrā)}{আক্রমণ করা (ākrômôn kôrā)}]{\bn{আক্রমণ} \bn{করা} (ākrômôn kôrā) — to attack}

\label{lesson:BN-C314-akromon-kora}



\begin{quote}
Before the new one: say the Bengali for an enemy, then the Bengali for a soldier.
\end{quote}

\subsection*{You'll want to know: \bn{আক্রমণ} \bn{করা}}
\textbf{\bn{আক্রমণ} \bn{করা}} — \emph{ākrômôn kôrā} — \textquotedblleft{}to attack\textquotedblright{}.

\begin{quote}
\bn{ভোরে} \bn{শত্রুরা} \bn{দুর্গ} \bn{আক্রমণ} \bn{করল।}
\end{quote}

\begin{quote}
\emph{bhore shôtrurā durgô ākrômôn kôrlô.}
\end{quote}

\begin{quote}
At dawn the enemy attacked the fort.
\end{quote}

Past (he/she …): \bn{সেই} \bn{সেনাবাহিনী} \bn{গ্রাম} \bn{আক্রমণ} \bn{করল}, \emph{sei senābāhinī grām ākrômôn kôrlô}.

\begin{cousinweb}
\textbf{\bn{আক্রমণ} \bn{করা}} (\emph{ākrômôn kôrā}) puts Sanskrit \emph{ākramaṇa}, 'a stepping upon, a mounting' (Monier-Williams), from \emph{ā-√kram}, 'to step on', before \bn{করা}; the sense 'an attack' is the later one.
\end{cousinweb}

\begin{grammarlens}[title={What you've built}]
One more word for telling what happened.
\end{grammarlens}

\subsection*{Your turn}
\begin{itemize}
\raggedright
  \item \emph{Say it:} \emph{ākrômôn kôrā}
  \item \emph{Say it:} \emph{ākrômôn kôrā}, once more
  \item \emph{Say it:} \emph{bhore shôtrurā durgô ākrômôn kôrlô.} — At dawn the enemy attacked the fort.
  \item \emph{Recall:} say the Bengali for an enemy, then the Bengali for a soldier, then say \emph{ākrômôn kôrā} again
\end{itemize}

\begin{tcolorbox}[breakable,colback=teal!4,colframe=teal!35!black,title={Before you move on}]
What does \bn{আক্রমণ} \bn{করা} mean? (\textquotedblleft{}To attack\textquotedblright{}.) Say it once more.
\end{tcolorbox}
\section[\texorpdfstring{\bn{হামলা} (hāmlā)}{হামলা (hāmlā)}]{\bn{হামলা} (hāmlā) — a raid; an assault}

\label{lesson:BN-C314-hamla}



\begin{quote}
Before the new one: say the Bengali for a soldier, then the Bengali for to attack.
\end{quote}

\subsection*{You'll want to know: \bn{হামলা}}
\textbf{\bn{হামলা}} — \emph{hāmlā} — \textquotedblleft{}a raid; an assault\textquotedblright{}.

\begin{quote}
\bn{রাতে} \bn{ডাকাতরা} \bn{গ্রামে} \bn{হামলা} \bn{করল।}
\end{quote}

\begin{quote}
\emph{rāte ḍākātrā grāme hāmlā kôrlô.}
\end{quote}

\begin{quote}
At night the robbers raided the village.
\end{quote}

\begin{cousinweb}
\textbf{\bn{হামলা}} (\emph{hāmlā}) is usually traced through Persian to Arabic \emph{ḥamla}, 'an attack, a charge'. \bn{হামলা} \bn{করা}, to raid.
\end{cousinweb}

\begin{grammarlens}[title={What you've built}]
One more word for telling what happened.
\end{grammarlens}

\subsection*{Your turn}
\begin{itemize}
\raggedright
  \item \emph{Say it:} \emph{hāmlā}
  \item \emph{Say it:} \emph{hāmlā}, once more
  \item \emph{Say it:} \emph{rāte ḍākātrā grāme hāmlā kôrlô.} — At night the robbers raided the village.
  \item \emph{Recall:} say the Bengali for a soldier, then the Bengali for to attack, then say \emph{hāmlā} again
\end{itemize}

\begin{tcolorbox}[breakable,colback=teal!4,colframe=teal!35!black,title={Before you move on}]
What does \bn{হামলা} mean? (\textquotedblleft{}A raid; an assault\textquotedblright{}.) Say it once more.
\end{tcolorbox}
\section[\texorpdfstring{\bn{দখল} \bn{করা} (dôkhôl kôrā)}{দখল করা (dôkhôl kôrā)}]{\bn{দখল} \bn{করা} (dôkhôl kôrā) — to occupy; to seize}

\label{lesson:BN-C314-dokhol-kora}



\begin{quote}
Before the new one: say the Bengali for to attack, then the Bengali for a raid.
\end{quote}

\subsection*{You'll want to know: \bn{দখল} \bn{করা}}
\textbf{\bn{দখল} \bn{করা}} — \emph{dôkhôl kôrā} — \textquotedblleft{}to occupy; to seize\textquotedblright{}.

\begin{quote}
\bn{শত্রুরা} \bn{তিন} \bn{দিনে} \bn{দুর্গটা} \bn{দখল} \bn{করল।}
\end{quote}

\begin{quote}
\emph{shôtrurā tin dine durgôṭā dôkhôl kôrlô.}
\end{quote}

\begin{quote}
The enemy took the fort in three days.
\end{quote}

Past (he/she …): \bn{তিনি} \bn{জমিটা} \bn{দখল} \bn{করলেন}, \emph{tini jômiṭā dôkhôl kôrlen}.

\begin{cousinweb}
\textbf{\bn{দখল} \bn{করা}} (\emph{dôkhôl kôrā}) puts \bn{দখল} — usually traced through Persian to Arabic \emph{dakhl}, 'entry' — before \bn{করা}: to get in and stay, to occupy.
\end{cousinweb}

\begin{grammarlens}[title={What you've built}]
One more word for telling what happened.
\end{grammarlens}

\subsection*{Your turn}
\begin{itemize}
\raggedright
  \item \emph{Say it:} \emph{dôkhôl kôrā}
  \item \emph{Say it:} \emph{dôkhôl kôrā}, once more
  \item \emph{Say it:} \emph{shôtrurā tin dine durgôṭā dôkhôl kôrlô.} — The enemy took the fort in three days.
  \item \emph{Recall:} say the Bengali for to attack, then the Bengali for a raid, then say \emph{dôkhôl kôrā} again
\end{itemize}

\begin{tcolorbox}[breakable,colback=teal!4,colframe=teal!35!black,title={Before you move on}]
What does \bn{দখল} \bn{করা} mean? (\textquotedblleft{}To occupy; to seize\textquotedblright{}.) Say it once more.
\end{tcolorbox}
